Las propiedades generales de la lisura que siguen son inmediatas.

\underline{Proposición 2.3.} Sean
\[
 \underline A\xrightarrow{\ \varphi\ }\underline B
 \xrightarrow{\ \psi\ }\underline C
\]
funtores de grupoides cofibrados sobre $\underline C_\Lambda$.

(a) Si $\varphi$ y $\psi$ son lisos, entonces también lo es
$\psi\circ\varphi$.

(b) Si $\psi\circ\varphi$ es liso y $\varphi$ es sobreyectivo sobre
las clases de isomorfía de objetos, entonces $\psi$ es liso.

%%%%%%%%%%  scan idx 63  |  folio 52  %%%%%%%%%%

\underline{Observación 2.4.} Sea $\underline A$ un grupoide cofibrado
sobre $\underline C_\Lambda$ que admite un objeto formal versal $\xi$.
Por 2.2(d), esto significa que
\[
 \tilde\xi:\underline h_R\longrightarrow\underline A
\]
es liso, donde $R=P(\xi)$. Por consiguiente, las afirmaciones
siguientes son equivalentes:

\begin{itemize}
\item $\underline A$ es formalmente liso sobre $\underline C_\Lambda$
      (cf. \sisi{1.13}{1.14}), es decir, $R$ es un álgebra de series formales sobre
      $\Lambda$;
\item el funtor compuesto
      \[
       \underline h_R\xrightarrow{\ \tilde\xi\ }\underline A
       \xrightarrow{\ P\ }\underline C_\Lambda
       \;(=\underline h_\Lambda)
      \]
      es liso;
\item $P:\underline A\longrightarrow\underline C_\Lambda$ es
      liso.
\end{itemize}
La última equivalencia se deduce de la Proposición 2.3(b).

Consideremos ahora una propiedad más fuerte que la semihomogeneidad.

\underline{Definición 2.5.} Un grupoide cofibrado $\underline A$ sobre
$\underline C_\Lambda$ se llama \underline{homogéneo} si todo diagrama
\[
\begin{tikzcd}[row sep=large,column sep=huge]
 & a' \arrow[d]\\
 b' \arrow[r] & a
\end{tikzcd}
\]
en el que $a'\longrightarrow a$ es sobreyectivo admite un producto
fibrado $a'\times_a b'$ en $\underline A$ tal que
\[
 P(a'\times_a b')=
 P(a')\times_{P(a)}P(b').
\]

\underline{Observación 2.6(a).} Sea $\underline A$ un grupoide
cofibrado sobre $\underline C_\Lambda$, y consideremos un diagrama
\[
(*)\qquad
\begin{tikzcd}[row sep=small,column sep=small]
 & a' \arrow[dl] \arrow[dr] & \\
 a_1 \arrow[dr] & & a_2 \arrow[dl]\\
 & a &
\end{tikzcd}
\qquad\text{en $\underline A$ situado sobre el diagrama}\qquad
\begin{tikzcd}[row sep=small,column sep=small]
 & R_1\times_R R_2 \arrow[dl,"{\pi_1}"'] \arrow[dr,"{\pi_2}"] & \\
 R_1 \arrow[dr] & & R_2 \arrow[dl]\\
 & R &
\end{tikzcd}
\]
en $\underline C_\Lambda$, donde los $\pi_i$ son las proyecciones.
Si existe un producto fibrado $a_1\times_a a_2$ en $\underline A$,
hay un único morfismo
\[
 \varphi:a'\longrightarrow a_1\times_a a_2
\]
compatible con las proyecciones. Si
\[
 P(a_1\times_a a_2)=R_1\times_R R_2,
\]
entonces $P(\varphi)$ es el morfismo identidad de $R_1\times_R R_2$;
por tanto, $\varphi$ es un $\underline C_\Lambda$-isomorfismo, es
decir, $a'=a_1\times_a a_2$. Así, si $\underline A$ es homogéneo,
todo diagrama $(*)$ define un producto fibrado $a_1\times_a a_2$
siempre que $R_1\longrightarrow R$ sea sobreyectivo. De ello se deduce
inmediatamente

%%%%%%%%%%  scan idx 64  |  folio 53  %%%%%%%%%%

que todo grupoide cofibrado homogéneo es, con mayor razón,
semihomogéneo.

\underline{Observación 2.6(b).} Consideremos un diagrama de categorías
y funtores
\[
\begin{tikzcd}[row sep=large,column sep=large]
 X \arrow[dr,"{\varphi}"'] & & Y \arrow[dl,"{\psi}"]\\
 & Z &
\end{tikzcd}
\]
El producto fibrado ordinario $X\times_ZY$ es la categoría cuyos
objetos son los pares $(x,y)$ que satisfacen
$\varphi(x)=\psi(y)$ y cuyos morfismos
$(x,y)\longrightarrow(x',y')$ son los pares
\[
 \alpha:x\longrightarrow x',\qquad
 \beta:y\longrightarrow y'
\]
tales que $\varphi(\alpha)=\psi(\beta)$. El producto fibrado
$X\times_Z^2Y$ en el sentido de las $2$-categorías es la categoría
cuyos objetos son las ternas $(x,y;\gamma)$, donde
\[
 \varphi(x)\xrightarrow{\ \gamma\ }\psi(y)
\]
es un isomorfismo. Un morfismo
$(x,y;\gamma)\longrightarrow(x',y';\gamma')$ es un par formado por
$\alpha:x\longrightarrow x'$ y $\beta:y\longrightarrow y'$ para el
cual el diagrama
\[
\begin{tikzcd}[row sep=huge,column sep=huge]
 \varphi(x) \arrow[r,"{\varphi(\alpha)}"] \arrow[d,"{\gamma}"']
 & \varphi(x') \arrow[d,"{\gamma'}"]\\
 \psi(y) \arrow[r,"{\psi(\beta)}"'] & \psi(y')
\end{tikzcd}
\]
es conmutativo.

Las nociones de grupoide cofibrado (semi)homogéneo y de funtor liso
pueden reformularse en términos de productos fibrados de categorías en
este sentido $2$-categórico. Sea
$P:\underline A\longrightarrow\underline C_\Lambda$ una categoría
cofibrada, y consideremos $R_1\times_RR_2$ en
$\underline C_\Lambda$. Elijamos funtores de imagen directa
\[
 (\pi_i)_*:\underline A(R_1\times_RR_2)
 \longrightarrow\underline A(R_i),\qquad i=1,2.
\]
Estos inducen un funtor
\[
 \underline A(R_1\times_RR_2)\longrightarrow
 \underline A(R_1)\times^2_{\underline A(R)}\underline A(R_2).
\]
Salvo equivalencia natural canónica, este no depende de los funtores
de imagen directa elegidos. El grupoide $\underline A$ es homogéneo si
y solo si este funtor es una equivalencia siempre que
$R_1\longrightarrow R$ sea sobreyectivo; es semihomogéneo si y solo si
\[
 [\underline A(R_1\times_RR_2)]\longrightarrow
 [\underline A(R_1)]
 \times_{[\underline A(R)]}[\underline A(R_2)]
\]

%%%%%%%%%%  scan idx 65  |  folio 54  %%%%%%%%%%

es sobreyectivo siempre que $R_1\longrightarrow R$ sea sobreyectivo,
y es biyectivo siempre que $R=k$ y $R_1=k[\varepsilon]$.

Sea ahora $\varphi:\underline A\longrightarrow\underline B$ un
funtor de grupoides cofibrados sobre $\underline C_\Lambda$. Para todo
morfismo $R'\longrightarrow R$ en $\underline C_\Lambda$,
$\varphi$ induce, salvo equivalencia natural canónica, un funtor
\[
 \underline A(R')\longrightarrow
 \underline A(R)\times^2_{\underline B(R)}\underline B(R').
\]
El funtor $\varphi$ es liso si y solo si este funtor es sobreyectivo
sobre las clases de isomorfía de objetos siempre que
$R'\longrightarrow R$ sea sobreyectivo.

\underline{Observación 2.6(c).} Si $\underline A$ es homogéneo,
entonces, para cualesquiera $k$-espacios vectoriales de dimensión
finita $V$ y $W$, el funtor
\[
 \underline A(k[V\times W])\longrightarrow
 \underline A(k[V])\times\underline A(k[W])
\]
es una equivalencia de categorías. Recordemos nuestro convenio
$\underline A(k)=\{e\}$. En particular,
\[
 \operatorname{Aut}_{\underline A(k[V\times W])}(1_{V\times W})
 \xrightarrow{\ \sim\ }
 \operatorname{Aut}_{\underline A(k[V])}(1_V)
 \times
 \operatorname{Aut}_{\underline A(k[W])}(1_W),
\]
donde $1_V$ es la imagen directa de $e$ por el único morfismo
$k\longrightarrow k[V]$ en $\underline C_\Lambda$. En consecuencia,
si $\underline A$ es homogéneo, tanto
$\operatorname{Aut}(1_\varepsilon)$ como
$[\underline A(k[\varepsilon])]$ están provistos de estructuras
canónicas de espacio vectorial sobre $k$.

\underline{Observación 2.6(d).} Para todo
$R\in\widehat{\underline C}_\Lambda$, el grupoide discreto cofibrado
$\underline h_R$ sobre $\underline C_\Lambda$ es homogéneo.

Si $\underline A$ es semihomogéneo, también lo es $[\underline A]$.
Sin embargo, la homogeneidad de $\underline A$ no implica en general
la homogeneidad de $[\underline A]$.

\underline{Proposición 2.7.} Sea $\underline A$ un grupoide cofibrado
homogéneo sobre $\underline C_\Lambda$. Las condiciones siguientes son
equivalentes.

(i) $[\underline A]$ es homogéneo; es decir,
\[
 [\underline A(R'\times_RS)]\longrightarrow
 [\underline A(R')]
 \times_{[\underline A(R)]}[\underline A(S)]
\]
es biyectivo para toda pequeña extensión $R'\longrightarrow R$.

(ii) Para toda pequeña prolongación $a'\longrightarrow a$ en
$\underline A$, el morfismo
\[
 \operatorname{Aut}_{\underline A(R')}(a')
 \longrightarrow
 \operatorname{Aut}_{\underline A(R)}(a)
\]

%%%%%%%%%%  scan idx 66  |  folio 55  %%%%%%%%%%

es sobreyectivo, donde $R'=P(a')$ y $R=P(a)$.

\underline{Demostración.} (i)$\Rightarrow$(ii). Sea
$f:R'\longrightarrow R$ una pequeña extensión en
$\underline C_\Lambda$, y sea $a'\longrightarrow a$ un $f$-morfismo
en $\underline A$. Dado
$\tau\in\operatorname{Aut}_{\underline A(R)}(a)$, consideremos
$a'\times_a a'$ y $a'\times_a a'_\tau$, donde
$a'_\tau\longrightarrow a$ es el compuesto
\[
 a'\longrightarrow a\xrightarrow{\ \tau\ }a.
\]
Como
\[
 [\underline A(R'\times_RR')]
 \longrightarrow
 [\underline A(R')]
 \times_{[\underline A(R)]}[\underline A(R')]
\]
es biyectivo, existe un isomorfismo
\[
 a'\times_a a'\xrightarrow{\ u\ }a'\times_a a'_\tau.
\]
La proyección sobre el segundo factor da entonces un isomorfismo
$\tau':a'\longrightarrow a'$ que levanta $\tau$.

(ii)$\Rightarrow$(i). Sean $a''_i$ $(i=1,2)$ objetos de
$\underline A(R'\times_RS)$ con la misma imagen por
\[
 [\underline A(R'\times_RS)]\longrightarrow
 [\underline A(R')]
 \times_{[\underline A(R)]}[\underline A(S)].
\]
Obtenemos un diagrama de flechas continuas
\[
\begin{tikzcd}[column sep=small,row sep=small]
 && a''_2 \arrow[ddll] \arrow[ddrr] &&\\
 && a''_1 \arrow[u,dashed] \arrow[dl] \arrow[dr] &&\\
 a'_2 \arrow[ddrr] & a'_1 \arrow[l,"{\sim}"] \arrow[dr]
 && b_1 \arrow[r,"{\tau}"] \arrow[dl] & b_2 \arrow[ddll]\\
 && a_1 \arrow[d,"{\bar\tau}"] &&\\
 && a_2 &&
\end{tikzcd}
\qquad\text{situado sobre}\qquad
\begin{tikzcd}[column sep=small,row sep=small]
 & R'\times_RS \arrow[dl] \arrow[dr] &\\
 R' \arrow[dr] && S \arrow[dl]\\
 & R &
\end{tikzcd}
\]
donde $\bar\tau:a_1\longrightarrow a_2$ es el isomorfismo inducido
por $\tau$, de modo que el cuadrado inferior derecho es conmutativo.
Bajo la condición (ii), el isomorfismo dado
$a'_1\xrightarrow{\sim}a'_2$ puede sustituirse de modo que también el
cuadrado izquierdo sea conmutativo. Como
\[
 a''_i=a'_i\times_{a_i}b_i\qquad(i=1,2)
\]
por la Observación 2.6(a), obtenemos un isomorfismo
$a''_1\xrightarrow{\sim}a''_2$. Por tanto,
\[
 [\underline A(R'\times_RS)]\longrightarrow
 [\underline A(R')]
 \times_{[\underline A(R)]}[\underline A(S)]
\]
es inyectivo y, por consiguiente, biyectivo, pues $\underline A$ es
homogéneo por hipótesis. \hfill$\square$
